\subsection{Shape-Aware Source Selection}\label{sec:sass}

This section presents our approach for shape-aware source selection.
\Cref{fig:running-example} presents the shapes of a running example and their \acp{LCG},
for a federation of two endpoints described by the shape index of \cref{lst:running-example-si},
so $e_1$ stores the users and $e_2$ stores the hobbies and clubs.


We describe each endpoint with the shapes it stores,
translated into \acp{LCG}, which are small \acp{KG} with constraints.
Each shape becomes an \ac{LCG} whose subject is the \ac{IRI} of the shape,
with the triples of its triple constraints, as for \texttt{ex:UserShape} in \cref{fig:running-example}.
The object of such a triple depends on the constraint on the object.
A reference to another shape becomes the \ac{IRI} of that shape,
which is the subject of its own \ac{LCG} when that shape is part of the shape index, so that the \acp{KG} are connected like the data they describe.
Allowed values are kept as they are,
and any other constraint becomes a fresh \ac{IRI}, such as \texttt{ex:l1}, standing for the terms satisfying it.
A shape with several alternatives can either be translated into a single \ac{KG}, in the \emph{coarse} \acsp{LCG},
or into one \ac{KG} per alternative, in the \emph{refined} \acsp{LCG},
as shown for \texttt{ex:UserShape} in \cref{fig:running-example}.
The coarse \ac{LCG} of a shape thus holds the union of the triples of all its alternatives.
Since no alternative forbids the triples of another, each triple of this union can occur in the data,
and only which triples occur together is lost from the \ac{KG}, while it remains in the constraints of the \ac{LCG}.\footnote{\Cref{def:shape} has no negation, so an alternative states which triples are allowed but never forbids a triple. Other criteria could build the coarse \acp{LCG}, as long as the \ac{LCG} of a set of alternatives remains looser than the \acp{LCG} of its subsets, accepting every triple they accept.}
Each shape is identified by an \ac{IRI}, and we write $\mathit{sh}(n)$ for the shape identified by the \ac{IRI} $n$.

\begin{definition}[\aclp{LCG} of a shape index]\label{def:summary}
Let $SI$ be a shape index.
Each triple constraint $\tau = (p, C_O, r)$ of a shape of $SI$ has the set of objects $\mathit{obj}(\tau) = \bigcup_{c \in C_O} \mathit{obj}(c)$, where
\[
  \mathit{obj}(c) =
  \begin{cases}
    \{ n_k \} & \text{if } c = @s_k \text{ with } s_k = \mathit{sh}(n_k) \in \mathrm{dom}(SI), \\
    \{ v \} \cup N(c) & \text{if } c \text{ is the single allowed value } v, \\
    \{ \iota_{C_O} \} \cup N(c) & \text{otherwise},
  \end{cases}
\]
where $\iota_{C_O}$ is a fresh \ac{IRI}, called the \emph{identifier} of $C_O$ and shared by every triple constraint with the same constraint $C_O$, that stands in the \ac{KG} for the terms satisfying $C_O$, while $C_O$ itself is kept outside the \ac{KG}, in the constraints of the \ac{LCG},
and $N(c) = \{ n_k \mid s_k = \mathit{sh}(n_k) \in \mathrm{dom}(SI) \text{ and } c \text{ can reach } s_k \}$ is the set of \acp{IRI} of the shapes that $c$ \emph{can reach},
which holds when $c$ admits an \ac{IRI} of $\mathit{ent}(s_k)$.\footnote{Blank nodes are left out, since they are scoped by the \ac{KG} of their endpoint.}
For an \ac{IRI} $n$ and a shape $s = (t, C_E, \mathit{TC})$, the \ac{LCG} $\pi(n, s) = (S, C)$ consists of the \ac{KG}
$S = \{ (n, p, o) \mid A \in \mathit{TC},\ \tau = (p, C_O, r) \in A,\ o \in \mathit{obj}(\tau) \}$
and the constraints $C = \{ n \mapsto s \} \cup \{ \iota_{C_O} \mapsto C_O \mid A \in \mathit{TC},\ \tau = (p, C_O, r) \in A,\ \iota_{C_O} \in \mathit{obj}(\tau) \}$.
A shape $s_i = \mathit{sh}(n_i) = (t_i, C_{E,i}, \mathit{TC}_i)$ is converted into $\sigma(s_i) = \{ \pi(n_i, s_i) \}$ in the \emph{coarse} \acsp{LCG},
and into $\sigma(s_i) = \{ \pi(n_i, (t_i, C_{E,i}, \{ A_j \})) \mid A_j \in \mathit{TC}_i \}$ in the \emph{refined} \acsp{LCG}.
The \acp{LCG} of $SI$ map each endpoint to the \acp{LCG} of the shapes it stores,
$\Psi(SI) = \{ e \mapsto \bigcup \{ \sigma(s_i) \mid s_i \in \mathrm{dom}(SI) \land e \in \mathit{IRI}_i \} \mid e \in E_{SI} \}$.
\end{definition}

A reference $@s_k$ to a shape of $SI$ produces the \ac{IRI} $n_k$ of $s_k$ as object, which is also the subject of the \acp{KG} of $s_k$,
so the \acp{KG} of two linked \acp{LCG} share a term even when their shapes are stored in different endpoints.

\begin{definition}[Link]\label{def:link}
Two \acp{LCG} are \emph{linked} when an \ac{IRI} in the \ac{KG} of one is the subject of the \ac{KG} of the other.
\end{definition}

The \acp{LCG} of $SI$ thus form a graph that is connected like the data they describe.

The options of a triple constraint thus determine the links of its \ac{LCG}.
A reference $@s_k$ links it exactly to the \acp{LCG} of $s_k$.
An option that admits \acp{IRI}, such as \texttt{IRI} for the objects of \texttt{foaf:knows} or an allowed \ac{IRI}, links it to the \acp{LCG} of every shape it can reach, since its objects can be entities of any of them.
This keeps every join of the data over the \acp{LCG}, but a query joining such an object with its own triples explores all these shapes, unless their targets or their constraints $C_E$ restrict their entities.
An option that admits only literals, such as \texttt{xsd:string}, links it to no other \ac{LCG}, since a literal is never the subject of a triple.

ASK queries can remove the links of an option that admits \acp{IRI} to the shapes that its objects never reach.
For such an option $c$ of a triple constraint $(p, C_O, r)$ of $s_i$ and each shape $s_k$ that $c$ can reach, an ASK query asks whether an object of $p$ in the endpoints of $s_i$ is the subject of a triple described by $s_k$ in the endpoints of $s_k$.
For a pair of endpoints $e_i \in \mathit{IRI}_i$ and $e_k \in \mathit{IRI}_k$, the query joins the objects of $p$ in $e_i$ with the subjects of $e_k$,
and restricts the triples of these subjects to the predicates that $s_k$ accepts, so that they are described by $s_k$ and not by another shape of $e_k$, as in \cref{lst:ask-reach}.
\begin{lstlisting}[caption={ASK query checking whether the \texttt{foaf:knows} objects of the users of \cref{fig:running-example}, if constrained only to be \acp{IRI}, reach \texttt{ex:ClubShape}.
The users are stored in \texttt{ex:e1} and the clubs in \texttt{ex:e2}.
The filter lists \texttt{ex:practices} and \texttt{ex:member}, the predicates that \texttt{ex:ClubShape} accepts, which excludes the hobbies stored in the same endpoint.
A positive answer keeps the link from \texttt{ex:UserShape} to \texttt{ex:ClubShape}, and a negative answer removes it.},label={lst:ask-reach}]
ASK {
  SERVICE ex:e1 { ?x foaf:knows ?o }
  SERVICE ex:e2 { ?o ?q ?v
    FILTER (?q IN (ex:practices, ex:member)) }
}
\end{lstlisting}
When the answer is negative, $n_k$ is removed from $N(c)$ for every triple constraint of $s_i$ with predicate $p$, since the query concerns every object of $p$,
while the identifier $\iota_{C_O}$ remains for the objects that are the subject of no triple.
The source selection remains sound, since a removed $n_k$ only matters for an object of $p$ that is the subject of a triple described by $s_k$, which the ASK query rules out.
Since the object and its triples can be stored in different endpoints, such a query spans several endpoints,
either as a federated ASK query with SERVICE clauses, as above, or by retrieving the objects of $p$ and asking the endpoints of $s_k$ about them,
at a cost that grows with the number of objects of $p$.
The pruned sets $N(c)$ can thus be materialized or cached, so that later queries reuse them without contacting the endpoints again, as long as the data does not change.
We extend the shape index vocabulary~\cite{tam2025} with \texttt{si:Reach}, which records for a shape and a predicate the shapes that its objects reach, together with the time of the ASK queries, so that the pruned links are published with the shape index.

In the coarse \acsp{LCG} of \cref{fig:running-example}, with $s_u = \mathit{sh}(\texttt{ex:UserShape})$, $\sigma(s_u)$ contains a single \ac{LCG},
whose \ac{KG} has the subject \texttt{ex:UserShape},
and whose constraints map \texttt{ex:l1}, \texttt{ex:l2}, and \texttt{ex:l3} to \texttt{LITERAL}, \texttt{xsd:string}, and \texttt{xsd:integer} with a maximum of 150.
In the refined \acsp{LCG}, $\sigma(s_u)$ contains one \ac{LCG} per alternative,
whose \acp{KG} share the subject \texttt{ex:UserShape} and whose constraints map it to the shape restricted to a different alternative.
The identifier \texttt{ex:l3} occurs in both \acp{LCG}, since \texttt{foaf:age} is outside the alternatives.
The \ac{KG} of \texttt{ex:HobbyShape} keeps the values \texttt{"sport"} and \texttt{"game"} of \texttt{ex:category}.
$\Psi(SI)$ maps $e_1$ to the \acp{LCG} of \texttt{ex:UserShape}, and $e_2$ to the \acp{LCG} of \texttt{ex:HobbyShape} and \texttt{ex:ClubShape}.

The refined \acsp{LCG} distinguish the alternatives of the shapes,
at the cost of a number of \acsp{LCG} that grows with the number of alternatives.
The coarse \acsp{LCG} can therefore be used alone when memory is limited or when the shapes have many alternatives,
at the cost of possibly selecting more sources.
More generally, the \acsp{LCG} could form several levels of refinement between the coarse and the refined \acsp{LCG}, which we leave for future work.

\ac{SASS} extends the result-aware source selection of FedUP~\cite{aimonierdavat:hal-04538238} with the information that only shapes provide,
namely which combinations of predicates an entity can have and which values its objects can take.
It keeps FedUP as its first stage, with $\sols$ evaluated over the coarse \acsp{LCG} instead of a HiBISCuS-type summary, which yields a unions-over-joins plan $\fmu(\Phi)$.
It then prunes the subplans of $\Phi$ that FedUP cannot rule out, in two shape-aware stages.
The second stage removes the subplans that have no solution over the refined \acsp{LCG}, because they combine predicates that no single alternative of a shape allows.
The third stage removes the subplans whose solutions require values that the shapes forbid,
which a constraint solver decides from the constraints of the matched terms together with the filters of the query,
whereas FedUP removes most filters on its summaries and replaces every literal by the same term.
Since the \acp{LCG} contain identifiers rather than the terms of the data, the query is first rewritten so that its evaluation over the \acp{LCG} only over-approximates its solutions,
and the information removed by this rewriting is restored as constraints in the third stage.
Each stage only prunes subplans, so the stages are ordered from the cheapest to the most expensive,
and the solver is only called on the subplans that survive the first two.

\lstinputlisting[caption={Shape index of the running example of \cref{fig:running-example}, serialized with the shape index vocabulary~\cite{tam2025}. Prefixes are omitted.},label={lst:running-example-si},basicstyle=\ttfamily\footnotesize]{code/example_shape_index.ttl}

\begin{figure}[t]
  \centering
  \lstset{basicstyle=\ttfamily\footnotesize, aboveskip=2pt, belowskip=2pt}
  \begin{minipage}[t]{0.48\textwidth}
    \centering\textbf{Shapes}
  \end{minipage}\hfill
  \begin{minipage}[t]{0.48\textwidth}
    \centering\textbf{Coarse \acsp{LCG}}
  \end{minipage}

  \begin{minipage}[t]{0.48\textwidth}
    \lstinputlisting[firstline=1, lastline=8]{code/example_shapes.shex}
  \end{minipage}\hfill
  \begin{minipage}[t]{0.48\textwidth}
    \lstinputlisting{code/example_summary_user.ttl}
  \end{minipage}

  \begin{minipage}[t]{0.48\textwidth}
    \lstinputlisting[firstline=9, lastline=12]{code/example_shapes.shex}
  \end{minipage}\hfill
  \begin{minipage}[t]{0.48\textwidth}
    \lstinputlisting{code/example_summary_hobby.ttl}
  \end{minipage}

  \begin{minipage}[t]{0.48\textwidth}
    \lstinputlisting[firstline=13, lastline=16]{code/example_shapes.shex}
  \end{minipage}\hfill
  \begin{minipage}[t]{0.48\textwidth}
    \lstinputlisting{code/example_summary_club.ttl}
  \end{minipage}

  \medskip
  \textbf{Refined \acsp{LCG} of \texttt{ex:UserShape}}

  \begin{minipage}[t]{0.48\textwidth}
    \lstinputlisting[firstline=1, lastline=4]{code/example_summary_user_refined.ttl}
  \end{minipage}\hfill
  \begin{minipage}[t]{0.48\textwidth}
    \lstinputlisting[firstline=6, lastline=10]{code/example_summary_user_refined.ttl}
  \end{minipage}
  \caption{Shapes of the running example and their coarse and refined \acp{LCG}.
  Prefixes are omitted.}
  \label{fig:running-example}
\end{figure}
